% main.tex — starting point for your own notes.
% Compile with:  latexmk main.tex        (or pdflatex twice)
% Full manual: README.md  ·  Macro reference: docs/MACROS.md
\documentclass[
  lang=english,     % english | italian
  accent=C09000,    % main colour (HTML hex), e.g. 2E7D32 green, CC7221 orange
  % circuits,       % uncomment to load circuitikz
  % qed=square,     % end-of-proof symbol: smiley (default) | square
]{opalith}

% ---- Document information (every field is optional) ------------------------
\title{Course Title}
\author{Your Name}
% \subtitle{Lecture notes}
% \headertitle{Short title}          % text in the odd-page header
% \professor{Prof. Jane Doe}
% \institution{University of Somewhere}
% \department{Department of Information Engineering}
% \degree{BSc in Computer Engineering}
% \coursecode{Course ID: IN01123456 \quad$\cdot$\quad 9 CFU}
% \academicyear{2026/2027}
% \logo[0.5\textwidth]{images/logo.png} % skipped if the file does not exist
% \date{}                             % empty = no date (default: today)

\begin{document}

\maketitle
\tableofcontents

\chapter{First chapter}

\section{First section}

Write your notes here. Key terms are \uline{underlined} and important words
are in \textbf{bold}.

\dfn{Name of the concept}{
  A definition goes inside a light box.
}

\thm{Name of the result}{
  Statement of the theorem.
}

\dem{}{
  The proof ends with the end-of-proof symbol.
}

\exer{Short title}{
  Text of the exercise.
}
\begin{Solution}
  Worked solution.
\end{Solution}

\end{document}
